\section{Institutional Background}\label{hist}

\subsection{Public procurement in S\~{a}o Paulo: the BEC platform}

Public procurement in the state of São Paulo is conducted through a decentralized system of purchasing agencies, each referred to as a Public Buyer Unit (PBU). Since 2007, the use of the state's electronic procurement platform, the \emph{Bolsa Eletrônica de Compras} (BEC), has been mandatory for PBUs when acquiring common goods and services. The BEC platform centralizes tender procedures and records detailed information on each procurement process. It plays a crucial role in São Paulo's procurement operations: for instance, in 2018 alone, approximately US\$1.7~billion worth of goods and services were procured via BEC. Over the period from 2009 onward, the BEC system handled more than R\$7.4~billion (approximately 34\% of all state government purchases) for the State Health Secretariat (SES/SP) and other agencies in São Paulo. Given the large scale and value of these tenders, ensuring fairness and preventing collusion in the procurement process is of paramount importance.

The BEC catalog contains a comprehensive listing of standardized items and services, classified by category. Nearly 55,000 purchase offers were issued through BEC during our study period, covering about 6,350 distinct items ranging from medical supplies to office equipment. Each purchase offer (PO) is a digital solicitation document that identifies one or more items to be procured, along with the required quantities and specifications. Each item within a PO is auctioned separately; thus a unique procurement transaction can be defined by a combination of a PO number and an item code (POI). The BEC records include rich information for each POI, including the reference (reserve) price set by the buyer, the number of participating firms, all bids submitted (identifying winners and losers), the auction format used, and the identities of the buyer (PBU) and auctioneer.

São Paulo's PBUs predominantly use two types of competitive tendering procedures available on BEC: (i)~a sealed-bid tender (\emph{convite}) and (ii)~a multi-round descending auction (\emph{pregão eletrônico}). In a sealed-bid tender, firms submit confidential bids by a specified deadline. After the submission period closes, all bids are opened and the contract is awarded to the qualified bidder with the lowest bid price. The sealed-bid format is typically used for purchases up to a value limit of R\$176,000 (approximately US\$35,000). In contrast, the descending auction (pregão) has no value limit and involves an open, interactive bidding phase. Initially, firms submit sealed proposals which are evaluated for qualification. Then, all qualified bids are announced (with the bidders' identities concealed) and an open auction phase commences. During this phase, which lasts for an initial 20 minutes, bidders can submit progressively lower bids, knowing the current lowest bid in real time. If a new bid is placed in the final minutes of this period (between 16 and 20 minutes), the auction is automatically extended by an additional 4 minutes. The auction concludes when no new bids arrive for 4~consecutive minutes. The winner is the bidder offering the lowest price at the end of this process, provided that price is below the government's preset reference price.

The key difference between the two formats lies in the opportunity for iterative bidding. The pregão format allows real-time bid competition and, after determining the lowest bid through the auction, a brief negotiation or confirmation step to potentially further reduce the price. In contrast, the convite is a single-shot process with no negotiation stage. From an oversight perspective, the open auction format provides more transparency during the bidding process (as all bids are observable in real time, albeit anonymously), whereas the sealed-bid format concentrates all transparency at the moment of opening. Both mechanisms are subject to anti-collusion regulations and audit processes, but their procedural differences can influence how collusion or favoritism might manifest. Our primary analysis focuses on pregão auctions, which account for the majority of BEC procurement by value and for which the close-margin regression-discontinuity framework is most naturally motivated; we report results on convite as a complementary robustness check.

\subsection{Antitrust enforcement: CADE and the conviction-anchored ground truth}

Brazil's federal competition authority, the \emph{Conselho Administrativo de Defesa Econ\^{o}mica} (CADE), is an autonomous agency within the Ministry of Justice responsible for investigating and adjudicating anticompetitive conduct. CADE's Tribunal renders administrative decisions (not criminal convictions in the judicial sense) that carry binding fines and behavioral remedies. Its investigative arm, the \emph{Superintend\^{e}ncia-Geral} (SG), conducts preliminary inquiries and administrative proceedings, often drawing on leniency agreements (\emph{Termos de Compromisso de Cessa\c{c}\~{a}o}, TCCs) in which cartel participants cooperate in exchange for reduced penalties. The typical enforcement pipeline proceeds from an initial complaint or leniency application, through a formal administrative proceeding, to a Tribunal decision that either convicts or acquits the named firms.

Our ground truth consists of CADE administrative proceedings in which the Tribunal found firms guilty of cartel conduct in product markets relevant to procurement in the state of S\~{a}o Paulo. We identified these cases through a systematic search of CADE's public case repository, restricting attention to proceedings with final Tribunal decisions (not pending or preliminary-stage cases) that specify a cartel period and list named firms with identifiable CNPJs (the Brazilian taxpayer identification number for legal entities). The resulting set contains seven distinct cartels operating in S\~{a}o Paulo procurement markets between 2004 and 2019, spanning sectors as diverse as pharmaceuticals, school meals, trash bags, solar water heaters, and metropolitan transit. Together, these cartels involve 35 convicted firms whose 8-digit CNPJ roots can be matched to the BEC bidding records.

The critical feature of this ground truth for our detection exercise is that CADE's case files delineate both the \emph{firms} and the \emph{product-market-time windows} over which each ring was found to have operated. This allows us to label individual BEC auctions as cartel-tainted---not merely firms as cartel-affiliated---by matching each convicted firm's CNPJ root and the CADE-specified conduct period to the corresponding BEC tenders by item code and year. The resulting conviction-anchored benchmark thus incorporates both the identity and timing dimensions of collusion, rather than treating all auctions involving a flagged firm as equally suspicious.
